% !TEX root = ../main.tex
\corrigesChapitre{2}{Méthodes physiques d'analyse}
\corrigesSouples

\begin{blocLivre}{cor-2-1}
\begin{corrige}{1}{Appliquer la loi de Beer-Lambert}
\demarche{Appliquer $A=\varepsilon\,\ell\,c$ en veillant aux unités, puis vérifier le domaine de validité.}
\begin{questions}
\item $A=\varepsilon\,\ell\,c=12\times1{,}00\times\num{5.0e-2}=\res{0{,}60}$.
\item $c=\dfrac{A}{\varepsilon\,\ell}=\dfrac{0{,}42}{12\times1{,}00}=\res{\qty{3,5e-2}{mol.L^{-1}}}$.
\item L'absorbance est proportionnelle à l'épaisseur traversée : elle double, $A=\res{1{,}2}$.
\item Une absorbance de 3,2 sort du domaine de validité de la loi : la solution est trop concentrée et
l'absorbance n'est plus proportionnelle à la concentration. Il faut \res{\text{diluer la solution}} avant la mesure.
\end{questions}
\end{corrige}
\end{blocLivre}
\begin{blocLivre}{cor-2-2}
\begin{corrige}{2}{Calculer une conductivité}
\demarche{Exprimer les concentrations en mol$\cdot$m$^{-3}$ avant d'appliquer la loi de Kohlrausch.}
\begin{questions}
\item $c=\qty{2,0}{mmol.L^{-1}}=\qty{2,0e-3}{mol}$ pour $\qty{1e-3}{m^3}$, soit $\res{c=\qty{2,0}{mol.m^{-3}}}$.
\item $\sigma=\left(\lambda(\ce{Na+})+\lambda(\ce{Cl-})\right)c=(5{,}01+7{,}63)\times2{,}0=\res{\qty{25,3}{mS.m^{-1}}}$,
soit $\res{\qty{0,253}{mS.cm^{-1}}}$ puisque 1 mS$\cdot$cm$^{-1}$ = 100 mS$\cdot$m$^{-1}$.
\item Pour le nitrate de potassium, $\sigma=(7{,}35+7{,}14)\times2{,}0=\qty{29,0}{mS.m^{-1}}$ : à concentration égale,
la solution de nitrate de potassium \res{\text{conduit mieux}}, car ses ions ont des conductivités molaires plus grandes.
\end{questions}
\end{corrige}
\end{blocLivre}
\begin{blocLivre}{cor-2-3}
\begin{corrige}{3}{La quantité de matière d'un gaz}
\demarche{Convertir dans les unités du système international, puis appliquer $n=pV/(RT)$.}
\begin{questions}
\item $T=20+273{,}15=\qty{293,15}{K}$ et $V=\qty{4,2e-3}{m^3}$, d'où
\[n=\frac{pV}{RT}=\frac{\num{1.05e5}\times\num{4.2e-3}}{8{,}314\times293{,}15}=\res{\qty{0,18}{mol}}
\qquad\text{et}\qquad m=n\,M=\res{\qty{0,72}{g}}.\]
\item $n=\dfrac{\num{2.00e7}\times\num{1.2e-2}}{8{,}314\times293{,}15}\approx\qty{98}{mol}$. Sous 200 bar, les
molécules sont proches les unes des autres : leurs interactions et leur volume propre ne sont plus négligeables,
\res{\text{le modèle du gaz parfait n'est qu'une approximation}}.
\end{questions}
\end{corrige}
\end{blocLivre}
\begin{blocLivre}{cor-2-4}
\begin{corrige}{4}{De la couleur au spectre}
\demarche{La couleur perçue est la complémentaire de la couleur absorbée.}
\begin{questions}
\item À 525 nm, la solution absorbe le \res{\text{vert}} ; sa couleur perçue est la complémentaire du vert,
le \res{\text{magenta}} (violet-rose).
\item Une solution jaune absorbe la couleur complémentaire du jaune, le bleu : la tartrazine absorbe
\res{\text{vers 430 nm}}, et non dans le rouge.
\item On règle le spectrophotomètre à la longueur d'onde du maximum d'absorption : environ 525 nm pour le
permanganate et 430 nm pour la tartrazine. La mesure y est la plus sensible.
\end{questions}
\end{corrige}
\end{blocLivre}
\begin{blocLivre}{cor-2-5}
\begin{corrige}{5}{Lire un spectre infrarouge}
\demarche{Repérer les bandes au-dessus de 1\,500 cm$^{-1}$, puis exploiter l'absence de certaines bandes.}
\begin{questions}
\item La bande large et forte vers 3\,300 cm$^{-1}$ correspond à la liaison \res{\text{O--H d'un alcool lié}} ;
les bandes vers 2\,900 cm$^{-1}$ correspondent aux liaisons C--H.
\item Aucune bande n'apparaît vers 1\,700 cm$^{-1}$ : la molécule ne contient pas de liaison C=O, ce qui élimine
l'éthanal et l'acide éthanoïque. De plus, la bande O--H n'a pas la largeur très étendue (2\,500 -- 3\,200 cm$^{-1}$)
d'un acide carboxylique. X est \res{\text{l'éthanol}} \ce{CH3-CH2-OH}.
\end{questions}
\schema[corrige]{02-ex05-ir}
\end{corrige}
\end{blocLivre}
\begin{blocLivre}{cor-2-6}
\begin{corrige}{6}{Doser un colorant alimentaire}
\demarche{Travailler à $\lambda_{\max}$, modéliser la droite d'étalonnage, puis tenir compte de la dilution.}
\begin{questions}
\item À 630 nm, le colorant absorbe le plus : la mesure y est la plus sensible et peu affectée par un léger
décalage de longueur d'onde.
\item Les points sont alignés sur une droite passant par l'origine : $A$ est proportionnelle à $c$, la loi de
Beer-Lambert est vérifiée. La modélisation donne $\res{k=\qty{0,1305}{L.\upmu mol^{-1}}}$
(soit $\varepsilon\approx\qty{1,3e5}{L.mol^{-1}.cm^{-1}}$).
\item Pour la solution diluée, $c=\dfrac{0{,}62}{0{,}1305}=\qty{4,75}{\upmu mol.L^{-1}}$ ; la boisson est deux fois
plus concentrée : $\res{c=\qty{9,50}{\upmu mol.L^{-1}}}$. En masse,
$c_\mathrm{m}=c\,M=\num{9.50e-6}\times792{,}8=\res{\qty{7,5}{mg.L^{-1}}}$.
\item La DJA vaut $6\times60=\qty{360}{mg}$ par jour, soit $\frac{360}{7{,}5}\approx\res{\qty{48}{L}}$ de boisson par jour :
une consommation normale reste très loin de la DJA.
\end{questions}
\schema[corrige]{02-ex06-cor}
\end{corrige}
\end{blocLivre}
\begin{blocLivre}{cor-2-7}
\begin{corrige}{7}{Contrôler un sérum physiologique}
\demarche{Diluer pour se placer dans la gamme d'étalonnage, modéliser, puis multiplier par le facteur de dilution.}
\begin{questions}
\item Prélever 5,0 mL de sérum à la pipette jaugée, les verser dans une fiole jaugée de 100,0 mL, compléter avec de
l'eau distillée jusqu'au trait de jauge et homogénéiser : le facteur de dilution vaut $\frac{100{,}0}{5{,}0}=20$.
Sans dilution (environ 154 mmol$\cdot$L$^{-1}$), la solution sortirait de la gamme d'étalonnage et du domaine de
validité de la loi de Kohlrausch.
\item La modélisation par une fonction linéaire donne $\res{k=\qty{0,1202}{mS.cm^{-1}}}$ par mmol$\cdot$L$^{-1}$.
\item Solution diluée : $c=\frac{0{,}928}{0{,}1202}=\qty{7,72}{mmol.L^{-1}}$ ; sérum : $c=20\times7{,}72=\res{\qty{154}{mmol.L^{-1}}}$.
En masse, $c_\mathrm{m}=\num{0.154}\times58{,}44=\res{\qty{9,0}{g.L^{-1}}}$ : l'indication de l'étiquette est vérifiée.
\end{questions}
\schema[corrige]{02-ex07-cor}
\end{corrige}
\end{blocLivre}
\begin{blocLivre}{cor-2-8}
\begin{corrige}{8}{Deux colorants dans un sirop}
\demarche{Exploiter d'abord la longueur d'onde où une seule espèce absorbe, puis l'additivité des absorbances.}
\begin{questions}
\item Le sirop absorbe dans le bleu (tartrazine) et dans le rouge-orangé (bleu patenté) : il transmet surtout la
lumière verte et paraît \res{\text{vert}}.
\item À 638 nm, seul le bleu patenté absorbe :
$c_\mathrm{B}=\dfrac{A_{638}}{\varepsilon\,\ell}=\dfrac{0{,}53}{\num{1.06e5}\times1{,}00}=\res{\qty{5,0e-6}{mol.L^{-1}}}$.
\item À 427 nm, $A_{427}=\varepsilon_\mathrm{J}\,\ell\,c_\mathrm{J}+\varepsilon_\mathrm{B}\,\ell\,c_\mathrm{B}$, d'où
\[c_\mathrm{J}=\frac{0{,}51-\num{3.0e3}\times\num{5.0e-6}}{\num{2.70e4}}=\res{\qty{1,83e-5}{mol.L^{-1}}}.\]
\item On trouverait $\frac{0{,}51}{\num{2.70e4}}=\qty{1,89e-5}{mol.L^{-1}}$, soit une erreur relative de \res{3~\%}.
\item Le sirop est dix fois plus concentré : $\res{c_\mathrm{B}=\qty{5,0e-5}{mol.L^{-1}}}$ et
$\res{c_\mathrm{J}=\qty{1,8e-4}{mol.L^{-1}}}$.
\end{questions}
\end{corrige}
\end{blocLivre}
\begin{blocLivre}{cor-2-9}
\begin{corrige}{9}{Trois isomères, trois spectres}
\demarche{Rechercher la bande C=O, la bande O--H et les bandes propres aux aldéhydes.}
\begin{questions}
\item Propanal \ce{CH3-CH2-CHO} : aldéhyde ; propanone \ce{CH3-CO-CH3} : cétone ;
prop-2-én-1-ol \ce{CH2=CH-CH2-OH} : alcool, avec une double liaison C=C.
\item Le spectre B présente une large bande O--H vers 3\,300 cm$^{-1}$ et une bande C=C vers 1\,645 cm$^{-1}$, sans
bande C=O : \res{\text{B est le prop-2-én-1-ol}}. Les spectres A et C présentent une forte bande C=O vers
1\,720 cm$^{-1}$ ; seul C présente en plus deux bandes fines vers 2\,820 et 2\,720 cm$^{-1}$, caractéristiques de la
liaison C--H d'un aldéhyde : \res{\text{C est le propanal}} et \res{\text{A la propanone}}.
\item Les bandes C--H de l'aldéhyde, entre 2\,650 et 2\,800 cm$^{-1}$, distinguent le propanal de la propanone.
\end{questions}
\end{corrige}
\end{blocLivre}
\begin{blocLivre}{cor-2-10}
\begin{corrige}{10}{Doser un comprimé effervescent par une mesure de pression}
\demarche{L'augmentation de pression, à volume et température constants, est due au seul gaz formé.}
\begin{questions}
\item $\Delta p=1\,352-1\,015=\qty{337}{hPa}=\qty{3,37e4}{Pa}$, d'où
\[n(\ce{CO2})=\frac{\Delta p\,V}{RT}=\frac{\num{3.37e4}\times\num{1.10e-3}}{8{,}314\times293{,}15}=\res{\qty{1,52e-2}{mol}}.\]
\item Une mole d'ions \ce{HCO3-} donne une mole de \ce{CO2} : $m=n\,M=\num{1.52e-2}\times84{,}0=\res{\qty{1,28}{g}}$.
L'écart relatif à l'étiquette vaut $\frac{1{,}30-1{,}28}{1{,}30}\approx\res{1{,}7~\%}$.
\item Le gaz resté dissous n'est pas compté : la quantité de \ce{CO2} mesurée, donc la masse calculée, est
\res{\text{sous-estimée}}. Cela explique une partie de l'écart.
\end{questions}
\end{corrige}
\end{blocLivre}
\begin{blocLivre}{cor-2-11}
\begin{corrige}{11}{La conductivité d'une eau minérale}
\demarche{Passer des concentrations en masse aux concentrations en quantité de matière, puis sommer les contributions.}
\begin{questions}
\item $c=c_\mathrm{m}/M$ : \ce{Ca^2+} 2,00 ; \ce{Mg^2+} 1,07 ; \ce{Na+} 0,283 ; \ce{K+} 0,026 ; \ce{HCO3-} 5,90 ;
\ce{SO4^2-} 0,146 ; \ce{Cl-} 0,282 ; \ce{NO3-} 0,061, en mmol$\cdot$L$^{-1}$.
\item Charges positives : $2\times2{,}00+2\times1{,}07+0{,}283+0{,}026=\qty{6,44}{mmol.L^{-1}}$ ; charges négatives :
$5{,}90+2\times0{,}146+0{,}282+0{,}061=\qty{6,53}{mmol.L^{-1}}$. Les deux valeurs sont \res{\text{très voisines}} :
l'eau est électriquement neutre, aux arrondis et aux ions non indiqués près.
\item Les concentrations en mmol$\cdot$L$^{-1}$ sont aussi des concentrations en mol$\cdot$m$^{-3}$ :
\[\sigma=\sum\lambda_i[X_i]=11{,}9\times2{,}00+10{,}6\times1{,}07+\dots+7{,}14\times0{,}061\approx\qty{67,9}{mS.m^{-1}}=\res{\qty{0,68}{mS.cm^{-1}}}.\]
\item La valeur mesurée vaut environ 85~\% de la valeur calculée. Les conductivités molaires tabulées correspondent
à des solutions très diluées ; ici, les ions, dont plusieurs portent deux charges, interagissent et conduisent
moins bien : \res{\text{la loi de Kohlrausch surestime la conductivité}}.
\end{questions}
\end{corrige}
\end{blocLivre}
\begin{blocLivre}{cor-2-12}
\begin{corrige}{12}{Contrôle qualité d'un antiseptique}
\demarche{Vérifier la sélectivité de la mesure, modéliser l'étalonnage, puis comparer à la tolérance de $\pm5$~\%.}
\begin{questions}
\item 525 nm correspond au maximum d'absorption de l'ion permanganate. L'ion hypochlorite n'absorbe que dans
l'ultraviolet, vers 292 nm : à 525 nm, \res{\text{seul l'ion permanganate absorbe}}.
\item Les points sont alignés avec l'origine ; la modélisation donne $\res{k=\qty{2,40e-3}{L.\upmu mol^{-1}}}$.
\item $c=\dfrac{0{,}150}{\num{2.40e-3}}=\qty{62,4}{\upmu mol.L^{-1}}$. Avec $M(\ce{KMnO4})=\qty{158,0}{g.mol^{-1}}$,
$c_\mathrm{m}=\num{62.4e-6}\times158{,}0=\qty{9,87e-3}{g.L^{-1}}$, soit \qty{0,99}{mg} pour 100 mL.
L'étiquette annonce 1,0 mg pour 100 mL : l'écart relatif vaut environ 1,3~\%, inférieur à 5~\%.
\res{\text{La solution est conforme}}.
\end{questions}
\end{corrige}
\end{blocLivre}
\begin{blocLivre}{cor-2-13}
\begin{corrige}{13}{Le gonflage d'un coussin de sécurité}
\demarche{Calculer la quantité de gaz nécessaire, puis remonter à la masse de réactif grâce à la stœchiométrie.}
\begin{questions}
\item La pression (1,2 bar) est faible et le diazote est loin de sa liquéfaction : \res{\text{le modèle est adapté}}.
\item $T=353{,}15$ K et $V=\qty{0,060}{m^3}$ :
$n(\ce{N2})=\dfrac{pV}{RT}=\dfrac{\num{1.20e5}\times\num{0.060}}{8{,}314\times353{,}15}=\res{\qty{2,45}{mol}}$.
\item Deux moles d'azoture donnent trois moles de diazote : $n(\ce{NaN3})=\frac23\,n(\ce{N2})=\qty{1,63}{mol}$.
Avec $M(\ce{NaN3})=22{,}99+3\times14{,}01=\qty{65,0}{g.mol^{-1}}$, $m=\res{\qty{106}{g}}$, soit une centaine de grammes.
\item À volume constant, $p$ est proportionnelle à $T$ : $p'=\num{1.20e5}\times\frac{293{,}15}{353{,}15}=\res{\qty{1,0e5}{Pa}}$.
La pression intérieure redevient égale à la pression atmosphérique : le coussin ne reste plus tendu et se vide
par ses évents, ce qui libère le passager.
\end{questions}
\end{corrige}
\end{blocLivre}
\begin{blocLivre}{cor-2-14}
\begin{corrige}{14}{De la menthe au laboratoire}
\demarche{Comparer les bandes O--H et C=O des deux spectres, puis les retrouver dans le spectre du produit brut.}
\begin{questions}
\item Menthol : groupe \res{\text{hydroxyle}} \ce{-OH}, famille des alcools ; menthone : groupe
\res{\text{carbonyle}} C=O, famille des cétones.
\item Menthol : \ce{C10H20O} ; menthone : \ce{C10H18O}. L'oxydation retire deux atomes d'hydrogène ; la chaîne
carbonée (cycle et ramifications) n'est \res{\text{pas modifiée}}.
\item Le spectre 2 présente une large bande O--H vers 3\,300 cm$^{-1}$ et aucune bande C=O : \res{\text{menthol}}.
Le spectre 1 présente une forte bande C=O vers 1\,710 cm$^{-1}$ sans bande O--H : \res{\text{menthone}}.
\item Le produit brut présente la bande C=O de la menthone, mais aussi une petite bande O--H vers 3\,400 cm$^{-1}$ :
il reste du menthol. La transformation est incomplète ou le produit n'est pas pur : \res{\text{il faut le purifier}},
puis vérifier la disparition de la bande O--H.
\end{questions}
\end{corrige}
\end{blocLivre}
\begin{blocLivre}{cor-2-15}
\begin{corrige}{15}{La salinité de l'eau de mer}
\demarche{Inverser le modèle d'étalonnage, tenir compte de la dilution, puis discuter les hypothèses.}
\begin{questions}
\item $a=\qty{0,1186}{mS.cm^{-1}}$ par mmol$\cdot$L$^{-1}$ et $b=\qty{0,012}{mS.cm^{-1}}$ sont la pente et l'ordonnée
à l'origine du modèle $\sigma=a\,c+b$. La seconde ligne affiche $c=\frac{\sigma-b}{a}=\res{\qty{5,90}{mmol.L^{-1}}}$,
concentration de l'eau de mer diluée.
\item L'eau de mer est cent fois plus concentrée : $c=\qty{0,590}{mol.L^{-1}}$, d'où une salinité
$c_\mathrm{m}=\num{0.590}\times58{,}44\approx\res{\qty{34,5}{g.L^{-1}}}$.
\item $\sigma=(5{,}01+7{,}63)\times590=\qty{7,46e3}{mS.m^{-1}}=\res{\qty{74,6}{mS.cm^{-1}}}$, valeur très éloignée de
la gamme d'étalonnage. À une telle concentration, la loi de Kohlrausch n'est plus valable : la dilution ramène
la mesure dans le domaine où le modèle a été établi.
\item Le résultat est une salinité « équivalente en chlorure de sodium » : les autres ions ont des conductivités
molaires et des masses molaires différentes. On trouve néanmoins une valeur très proche des 35 g$\cdot$kg$^{-1}$
attendus : \res{\text{l'estimation est pertinente}}, avec une précision de quelques pourcents.
\end{questions}
\schema[corrige]{02-ex15-cor}
\end{corrige}
\end{blocLivre}
